\section{Introduction}
\label{sec:intro}
Autonomous cyber-physical systems such as cars, drones and industrial robots run a processing pipeline that typically consists of sensing, perception, planning and control. The execution path from a sensor reading to the resulting actuation must complete within an end-to-end deadline for the system to be considered functional, safe, or both. Some steps of the pipeline require hardware acceleration and routinely offload work to GPUs. Such systems are usually built on a middleware, a software layer between the application and the operating system that provides communication and composition, and thereby raises the level of abstraction at which components are written. This is what lets the hundreds of components making up the pipeline be implemented with reasonable effort.

ROS~2~\cite{ros2-standard} is one such middleware, widely used in industry and academia. It abstracts an application into three levels: the callback is the unit of computation; a node owns the interfaces through which its callbacks interact with others; and an executor dispatches the ready callbacks of the nodes assigned to it on one or more threads of an OS process.

This creates the difficulty that hinders analysis, verification and validation of such systems: three timelines run at once, and none of them sees the other two (Fig.~\ref{fig:timelines}). From the point of view of the OS, threads are scheduled onto cores; from that of a ROS~2 executor, ready callbacks are dispatched onto its threads; from that of the GPU, offloaded work is scheduled onto its engines. When a callback runs, on which thread, core or GPU engine, and how long it takes is decided on all three: in Fig.~\ref{fig:timelines}, callback 2 is a single execution to the executor, two stretches on a core to the OS, and four pieces of queued work to the GPU, and the time it spends waiting---for the GPU behind another process's copy, then for a core after being preempted---appears on no single timeline. And the three timelines are only where the work runs; what work there is, and in what order, comes from the application's architecture---which callbacks exist, what triggers each, and how one's output becomes another's input---so that, too, has to be known.

Analysing, verifying or validating such a system therefore needs information from all three timelines and from the architecture that drives them: which callbacks exist and how they map to executors and threads, when each ran, what GPU work it submitted, and how callbacks interact. Timing analyses of ROS~2~\cite{casini-19-reservationbased,arafat-2026-mtexecutor-cbgroups,teper-2024-end2endmtexecutoranalysisoptimization} take this information as given, and in practice it is rarely written down or current; the tools that extract it from a running system each cover a part of it, and none the GPU (Section~\ref{sec:goal}).

In this work, we present FISH\footnote{GitHub: \url{https://github.com/FatihTasyaran/FISH_INTERFERE}}\textsuperscript{,}\footnote{Documentation: \url{https://fish-interfere.readthedocs.io}}, a toolset that extracts this information automatically from a traced run of an unmodified ROS~2 application: it is transparent to the application, applies to any application, and records all three timelines on one clock, so that each can be analysed on its own and all of them together. The model is based on actual execution traces. Therefore, it does not require expert knowledge about the application that needs to be modeled, nor does it rely on potentially incomplete or outdated documentation. \textbf{Our contributions are:} \textbf{(I)}~instrumentation that discovers the deployed architecture, records every callback execution with its thread---all five kinds the executor dispatches, for rclcpp and rclpy---attributes GPU work to the callback that submitted it, and recovers how callbacks interact; \textbf{(II)}~an acquisition procedure that traces any launched application unmodified, GPU processes included, and separates its phases; \textbf{(III)}~extraction of a hierarchical model that captures each callback's GPU behaviour per invocation and can be read at any granularity, shown feasible on Autoware, an industrial-size application. Section~\ref{sec:goal} states what is extracted and why it is hard, Section~\ref{sec:method} how FISH does it, Section~\ref{sec:findings} what it yields on Autoware.
